% Einheit 1 von 5 – Fläche zwischen Graph und x-Achse (bank/flaecheninhalt-durch-integration/e1.jsonl, zusammenbau v0.1)
%% TODO Zeitmarke Einheit 1: Marken-Zeile nicht lesbar
%% TODO Zweigzeile Teil 1: was hier gelernt wird (Satz in Schülersprache aus dem Einheitstitel) – Einheit 1
\einheitenkopf[e1]{Einheit 1 von 5 · Fläche zwischen Graph und x-Achse}
\zweigzeile{Abitur GK}
\setcounter{aufgabe}{8}

%% TODO Titel: Ich-kann-Satz fehlt in der Bank (Platzhalter: Kettenname) – Nr. 9
%% TODO Anweisung: ein Satz über der ersten Teilaufgabe fehlt in der Bank – Nr. 9
\begin{aufgabe}{Fläche zwischen Graph und x-Achse}
%% TODO \rechenplatz{n}: Schrittzahl des Grundfalls fehlt in der Bank (nur bei fester Schreibform, 2.3 b)
\begin{teile}
\teil Fläche zwischen dem Graphen von $f$ und der x-Achse von $x = -2$ bis $x = 2$ – wo liegt sie? \kreuz{ganz oberhalb der x-Achse}\\ \kreuz{ganz unterhalb der x-Achse}\\ \kreuz{auf beiden Seiten der x-Achse}

\begin{ksys}[xmin=-3,xmax=4,ymin=-5,ymax=5,ablesen]\funktion{0.5*\x^2+1}{f}\end{ksys}
\teil $f(x) = x^2 - 4x$ – Fläche, die der Graph mit der x-Achse einschließt?
\teil $f(x) = x^3 - 6x^2 + 11x - 6$ hat drei ganzzahlige Nullstellen – Fläche zwischen Graph und x-Achse zwischen den beiden kleinsten?
\teil $f(x) = -x^2 + 2x + 8$ – Fläche zwischen Graph, x-Achse und y-Achse im ersten Quadranten?
\teil $f(x) = x^3 - 9x$ – Gesamtfläche, die der Graph mit der x-Achse einschließt?
\teil Eine Giebelwand wird oben vom Graphen von $f(x) = -0{,}5x^2 + 8$ und unten von der x-Achse begrenzt; eine Längeneinheit entspricht $0{,}5$ m. Die Tür hat $2\,\text{m}^2$ – zu streichende Wandfläche?
\teil Die Abbildung zeigt den Graphen von $f$. Für den Inhalt der Fläche zwischen Graph und x-Achse von $x = 0$ bis $x = 6$ hat jemand $12$ erhalten. Entscheide mithilfe der Abbildung, ob das stimmen kann. \hfill (Abitur 2024 LK)

\begin{ksys}[xmin=-1,xmax=7,ymin=-5,ymax=2,ablesen]\funktion{0.5*\x^2-3*\x}{f}\end{ksys}
\end{teile}
\end{aufgabe}

%% TODO Titel: Ich-kann-Satz fehlt in der Bank (Platzhalter: Kettenname) – Nr. 10
%% TODO Anweisung: ein Satz über der ersten Teilaufgabe fehlt in der Bank – Nr. 10
\begin{aufgabe}{Fläche zwischen Graph und x-Achse – Fehler finden, Begründen, Anwendung}
\begin{teile}
\teil Fläche zwischen dem Graphen von $f(x) = x^2 - 2x$ und der x-Achse – wo ist der Fehler? \rechnung{F(x) &= \tfrac{1}{3}x^3 - x^2 \\ A &= F(2) - F(0) = \tfrac{8}{3} - 4 = -\tfrac{4}{3}}
\teil Warum nimmt man bei einer Fläche unterhalb der x-Achse den Betrag des Integrals?
\teil Eine Rasenfläche wird von einem Weg (x-Achse), einer Mauer (y-Achse), einer Hecke (Gerade $x = 5$) und einem Bach (Graph von $f(x) = 0{,}1x^2 + 2$) begrenzt; eine Längeneinheit entspricht $10$ m. Ein Sack Rasensamen reicht für $250\,\text{m}^2$ – wie viele Säcke?
\end{teile}
\end{aufgabe}
